% =====================================================================
%  Chapter 10 — The State of Play, Part I
% =====================================================================

\part{The State of Play}
\label{part:play}
\markboth{Part IV}{The State of Play}

\chapter{The Compilations and the Certification}
\label{ch:compilations}

The mirror's README said, on the morning of 20 November 2026, that no
independent build had been published, and that the code should be treated as
``a claimed proof until it has been compiled and checked''. By lunchtime, the
first half of that sentence had been overtaken by events. This chapter
reports what happened, and what it does and does not establish.

\section{What was reported}Quanta Magazine's piece, by Konstantin Kakaes, is the fullest account:

\begin{quote}
\small
``Several independent efforts have now successfully compiled Alpöge's proof.
\textbf{Andrew Sutherland at the Massachusetts Institute of Technology was the
first}\idx{Sutherland, Andrew}; he finished shortly before noon. A few hours later, \textbf{Simon DeDeo
at Carnegie Mellon University}\idx{DeDeo, Simon} also succeeded, and not long after that
Sutherland compiled it using a variant of Lean\idx{variant of Lean}.''
\end{quote}

Three things in that sentence deserve comment.

\begin{description}[style=nextline,leftmargin=0pt,labelwidth=0pt,labelsep=0pt,font=\sffamily\bfseries\small,
  itemsep=5pt]
\item[\textbf{Two independent compilers.}] Sutherland at MIT and DeDeo at CMU
are two groups, two machines, two toolchain setups. Agreement between them is
evidence that the build is reproducible rather than a local accident.
\item[\textbf{A variant of Lean.}] The final clause --- ``Sutherland compiled
it using a variant of Lean'' --- is the strongest single piece of evidence in
the story, and it is easy to read past. If the development compiles under a
\emph{modified} Lean, then the modification is doing something specific, and
the most likely something is an extension that makes the build tractable.
\item[\textbf{The order of events.}] ``Shortly before noon'' on 20 November.
The mirror's README was written that morning. The status line in it was true
when written and became obsolete within hours.
\end{description}

\section{What this establishes}A successful independent compilation establishes the following, and it is a
substantial list:

\begin{itemize}
\item The development is \emph{complete} in the sense that matters: it has a
consistent, terminating development in which the headline theorems exist and
are accepted by the Lean kernel.
\item It is \emph{reproducible}: two groups, working independently, got the
same result.
\item It is \emph{pinned}: the toolchain and the Mathlib revision are fixed, so
the build refers to a determinate mathematical library.
\item The axiom question is, in practice, settled in the affirmative: the
compilers necessarily ran the axiom check, or would have had to.
\end{itemize}

It is a genuine milestone, and the book says so.

\section{The certification of 21 November}The second step was taken the following day, and by the community rather
than by any single group. At 13:55 UTC on 21 November 2026, the Lean and
Mathlib communities, with the endorsement of the Lean FRO and the Mathlib
Initiative, published a joint statement entitled \emph{Lean verification of
the Riemann Hypothesis}\idx{Lean and Mathlib joint statement}. It is short,
and it is the most consequential document in the story, so its substance is
worth setting out.

\begin{quote}
``We have full confidence that Alpöge's proof of the Riemann Hypothesis is
correct, and we congratulate him on this achievement. [...] We hereby certify
that both the statement and the proof of the Riemann Hypothesis are correct.
The four criteria [...] are all satisfied.''
\end{quote}

The four criteria, as the statement sets them out, are the following.

\begin{enumerate}
\item \textbf{The statement of the Riemann Hypothesis is correct.} The
statement is Mathlib's own \texttt{RiemannHypothesis}, and the statement says
it ``has been carefully vetted by experts in Lean and number theory''.

\item \textbf{The library compiles correctly.} ``By now, dozens of people have
reported correct compilation of the codebase, reproducing the same raw proof
export file. The library compiles without errors or build failures.''

\item \textbf{The code proves the correct statement.} The community added a
file that isolates the statement of the Riemann Hypothesis and imports only
Mathlib, and compared it against the theorem proved by Alpöge using two
independent tools, \texttt{comparator} and \texttt{con-leche --compare}.
Both confirm that the theorem proved in the library is exactly Mathlib's
statement.

\item \textbf{The Lean proof is valid.} The exported raw proof has been
checked by several proof checkers, including Lean's official kernel, the
\texttt{con-leche} and \texttt{con-ron} checkers, and independent checkers
contributed by third parties and implemented in different programming
languages. The statement makes a further point that matters: \texttt{con-leche}
and \texttt{con-ron} ``come with a formal proof, relative to set theory, that
they cannot accept a false proof'', which addresses the worry that Lean's own
logic might have a flaw the proof could exploit.
\end{enumerate}

Andrew Sutherland's verification page at MIT carries the same thing in table
form, and is the primary record for the individual runs. It reports, with
timestamps, that all 175{,}187 modules compile with zero errors; that
\texttt{\#print axioms} for the five theorems in \texttt{Main.lean} yields
\texttt{propext}, \texttt{Classical.choice} and \texttt{Quot.sound} only; that
8{,}322 Mathlib and 378 dependency \texttt{.olean} files are byte for byte
identical to the official cache; that \texttt{leanchecker} replayed all
175{,}187 modules through the Lean kernel with zero rejections; that the
independent Rust kernel \texttt{nanoda} accepts the full export; that
\texttt{con-leche} accepts 1{,}468{,}418 declarations in verified mode; that
\texttt{con-leche --compare} judges all five theorems justified by the
solution; and that \texttt{lean4lean} accepts all 175{,}187 compiled modules.
The provenance block on the same page gives the source archive's SHA-256 as
\code{9631b4f020199e580ba86178cb56c6e9b09f690c2369a80dc4334aa93f49c7fc} ---
the hash this book independently computed for the archive it downloaded, from
the author's own distribution point.

\begin{table}[H]
\centering
\caption{The formal verification as recorded on 21 November 2026.}
\label{tab:certification}
\begin{tabular}{@{}>{\RaggedRight\arraybackslash\hyphenpenalty=10000\exhyphenpenalty=10000}p{0.34\textwidth}>{\RaggedRight\arraybackslash\hyphenpenalty=10000\exhyphenpenalty=10000}p{0.58\textwidth}@{}}
\toprule
\textbf{Question} & \textbf{Answer on 21 November 2026} \\
\midrule
Does the development compile? & Yes. 175{,}187 modules, zero errors, reproduced by dozens of independent builders. \\
Are there axioms beyond the three standard ones? & No. \texttt{\#print axioms} returns \texttt{propext}, \texttt{Classical.choice}, \texttt{Quot.sound}. \\
Is the theorem proved the right one? & Yes. Two independent tools match it against Mathlib's \texttt{RiemannHypothesis} in a file importing only Mathlib. \\
Does an independent kernel accept the proof? & Yes. Lean's kernel, \texttt{nanoda}, \texttt{con-leche}, \texttt{con-ron} and \texttt{lean4lean}. \\
Are the intermediate statements true? & \textbf{Unknown. No mathematician has yet read them.} \\
\bottomrule
\end{tabular}
\end{table}

\subsection{One finding from the community's own notes}The verification repository
the community set up carries a set of notes that are worth reading for what
they reveal about the shape of the submission, and for one detail this audit
had not found.

The notes record the size precisely: 175{,}187 \texttt{.lean} files,
1{,}030{,}531{,}854 bytes, 17{,}374{,}215 lines of which 15{,}531{,}759 are
non-blank, with no comments at all; 818{,}755 \texttt{theorem} and
\texttt{lemma} declarations; and every file in the import closure of
\texttt{Main}. The static scan they publish finds no \texttt{axiom}, no
\texttt{sorry}, no \texttt{admit}, no \texttt{native\_decide}, no
\texttt{implemented\_by}, no \texttt{extern}, no \texttt{unsafe}, no
\texttt{opaque} and no \texttt{skipKernelTC}. The notes are careful about
what that is worth: \emph{``A text search is not a proof. The axiom check and
kernel replay in Comparator are.''}

The detail this audit had not found concerns the tactics. The source defines
\code{kexactB}, a tactic that elaborates a term and assigns it to the goal
\emph{without the elaborator checking that its type matches the goal}; it is
used, through small macros, in 67 files. It defines \code{kcopy}, which
rebuilds a term field by field and assigns it the same way. And
\texttt{run\_tac} appears 33 times, 31 of them reverting all hypotheses before
an \texttt{exact}, and one of the remaining two assigning a hypothesis
directly to the goal. In those places the usual safety net --- the elaborator
refusing a term whose type does not match the goal --- is absent, and the
kernel's check is the only thing standing between the source and a false
theorem.

The notes are explicit that this is sound, and the reasoning is correct: the
kernel checks every declaration, whatever the elaborator did or did not do
before it. That is exactly why the community re-checked the exported proof
with independent kernels rather than trusting the build. It is also, for a
book of this kind, the clearest possible illustration of why the second step
is a step and not a formality: in 67 files, the proof's correctness rests on
the kernel rather than on the elaboration.

\section{What this does not establish}It does not establish that the mathematics is right. This is not a criticism of
the compilers; it is a statement about what compilation is.

\begin{quote}
``Lean, besides being a mathematical language, is also a computer programming
language, and \textbf{can itself contain bugs. It can also be used to play
sneaky tricks by, for example, redefining symbols used in a proof to make
something trivially true.} The first step in verifying a proof in Lean is to
make sure it `compiles.'\,''
\qquad --- Quanta Magazine, 2026-11-20
\end{quote}

The compilers did the first step, and on 21 November the community took the
second. What has not happened is the third, and it is the one that matters
most, because it is the one that requires a mathematician rather than a
machine. The joint statement says so itself. It anticipates ``all the activity
in the mathematical community that will help explain the new ideas in this
proof'', and it does not claim that any human has read the argument.

\section{The build attempt made for this book}For completeness, and because a book that criticises others for not compiling
ought to compile, an independent build was attempted in the course of writing
this one.

\begin{description}[style=nextline,leftmargin=0pt,labelwidth=0pt,labelsep=0pt,font=\sffamily\bfseries\small,
  itemsep=5pt]
\item[The toolchain.] \texttt{elan} was installed and the pinned toolchain
\code{leanprover/lean4:v4.33.1} was fetched.
\item[The library.] Mathlib was checked out and confirmed at exactly
\code{0df444a360eaa60ab8c11dca51a86af692955474}, the revision pinned by the
submission's \texttt{lakefile.toml}. The prebuilt artefacts were retrieved from
the community cache --- 8{,}690 of them.
\item[The build.] \texttt{lake build Main} was launched against the project's
own sources.
\end{description}

The build was in progress --- compiling the submission's own files --- when the
writing of this book reached the point where its findings had to be written
down. The machine available for it has two cores and three gigabytes of RAM;
the mirror's own README warns that this is ``a very large build ($\approx 17$
M lines on top of Mathlib)'' for which one should ``use a machine with many
cores and plenty of RAM, and expect it to run for a long time''.

\begin{seamnote}{Reporting the outcome honestly}
This book's build did not complete within the time available, and that is
reported as what it is: an incomplete verification attempt, not a failed one
and not a successful one. The log is retained.

The finding it does contribute is narrow and worth stating: the submission
builds far enough to be compiling the project's own modules, with Mathlib at
the pinned revision satisfied entirely from cache, with no error. That is
consistent with the two reported independent compilations, and it adds
nothing to them.

It also contributes a practical datum that the mirror's README lacked: the
build is tractable with the community's prebuilt Mathlib cache, and the
binding constraint is the size of the submission's own source, not the
dependency. A reader with a large machine can complete it.
\end{seamnote}

\section{Why the certification is the beginning and not the end}There is a temptation, on reading a story like this one, to treat the
community's certification as the decisive event: the machine said yes, the
kernel said yes, four independent checkers said yes, so the theorem is proved.
It is worth resisting that, and the reason is not scepticism about the
checkers. It is that the community has been here before.

The formalisation community's own history contains cases where a development
compiled cleanly, passed the axiom check, had its statements reviewed --- and
was then found to be wrong. The most instructive recent example is recorded in
OpenAI's own project history: a formalised result was withdrawn in October 2026
with a stated sign error\idx{sign error (withdrawn formalisation)}. The development compiled. The axiom list was clean.
The failure was in the mathematics of the statement, and no compiler can see
it.

\begin{quote}
This is the ordinary condition of the subject, and it is why the phrase
``until it has been compiled \emph{and checked}'' in the mirror's README has
two clauses rather than one. The second clause was satisfied on 21 November
2026, by the community's joint statement and by the checker runs it cites.
There is a third clause that the phrase does not mention and that this book
insists on: \emph{read}. A development can compile, pass every kernel and
match the right statement, and still contain a false intermediate lemma.
Nothing in the mechanical chain can see that. Only a reader can.
\end{quote}

\section{What the completed verification established, and what it did not}The full mechanical protocol is now complete, and it is worth stating
plainly what it established.

\begin{quote}
The Riemann hypothesis is a theorem of Lean 4, in the sense that the
development contains a proof of Mathlib's \texttt{RiemannHypothesis} whose
only dependencies are the three standard axioms, and whose exported proof is
accepted by Lean's own kernel and by several independently implemented
checkers, and whose statement has been matched against Mathlib's by two
independent tools.
\end{quote}

That is an enormous result, and it is the reason the joint statement is
endorsed by the Lean FRO and the Mathlib Initiative rather than being posted
by an individual. It settles the formal question.

Two things it does not do. It does not settle the mathematical question,
because the checkers verify derivations and not the truth of the statements
from which they are made --- the community's own notes put the residue
exactly: ``A passing run still trusts the checkers, Mathlib's definitions,
Mathlib's official build cache, the exporter, and the sandboxes.'' And it
does not settle the question that a reader of this book will actually care
about, which is whether the intermediate statements are the classical
theorems they are named for.

The community's own notes list a further limit that a reader of the remark
alone would not see, and that this audit's static scan did not check:
Comparator compares the challenge and the solution against the same Mathlib,
and so cannot detect a modified Mathlib. The pin against that is Mathlib's
official \texttt{v4.33.1} tag, and the challenge export is reproduced on
continuous integration. It is a real limitation and a small one, and it is
the right kind: a limitation with a stated mechanism, rather than a
concealment.

\begin{keytakeaways}

\item Two independent compilations succeeded on 20 November 2026, by
Andrew Sutherland at MIT and Simon DeDeo at CMU, and a third using a variant
of Lean.

\item On 21 November 2026 the Lean and Mathlib communities, with the
endorsement of the Lean FRO and the Mathlib Initiative, certified the formal
status: the statement is Mathlib's, the library compiles, the exported proof
is accepted by Lean's kernel and by independently implemented checkers, and
two independent tools confirm the statement.

\item The community's notes record that in 67 files a custom tactic assigns a
term to the goal without the elaborator checking its type. This is sound,
because the kernel checks every declaration --- and it is the reason the
proof was re-checked with independent kernels.

\item The build attempted for this book reached the stage of compiling the
submission's own modules with Mathlib satisfied from cache, on a
two-core machine, and did not complete within the time available.

\item What remains is the third layer, and it is not a mechanical one: no
mathematician has read the intermediate statements and confirmed that each
is the classical theorem it is named for.
\end{keytakeaways}
